\documentclass[reqno]{amsart}
\usepackage[T1]{fontenc}
\usepackage{amsmath,amssymb,amsthm}
\usepackage{xr-hyper}
\usepackage[hidelinks]{hyperref}
\externaldocument[tw-]{20261126T031000--twists-with-few-congruences__workbook}
\externaldocument[w2-]{20261126T060600--lengths-as-values-of-polynomials__workbook}

% Chapter number; \ref{chap} gives it.
\newcounter{chap}
% Chapter numbers, generated from 20261125T012000--grh-workbook__guide.org
% by scripts/make-chapters.py.  Do not edit: change the order of the
% index and rerun the script (./preview-site.sh and CI do this).
%
% Usage in a sheet, after \newcounter{chap}:
%   % Chapter numbers, generated from 20261125T012000--grh-workbook__guide.org
% by scripts/make-chapters.py.  Do not edit: change the order of the
% index and rerun the script (./preview-site.sh and CI do this).
%
% Usage in a sheet, after \newcounter{chap}:
%   % Chapter numbers, generated from 20261125T012000--grh-workbook__guide.org
% by scripts/make-chapters.py.  Do not edit: change the order of the
% index and rerun the script (./preview-site.sh and CI do this).
%
% Usage in a sheet, after \newcounter{chap}:
%   \input{chapters}\chapnumber{<sheet ID>}
\makeatletter
\providecommand\chapters@set[2]{\@namedef{chapters@#1}{#2}}
\chapters@set{20261125T001000}{1}% zeros-to-ramp-sums
\chapters@set{20261125T001100}{2}% inert-primes-and-estimate-b
\chapters@set{20261125T005000}{3}% galois-fourier-and-averaging
\chapters@set{20261125T005500}{4}% columns-from-a-sphere
\chapters@set{20261125T000000}{5}% slope-method
\chapters@set{20261126T030100}{6}% a-twist-for-each-prime
\chapters@set{20261125T003000}{7}% the-chain-and-the-square-root
\chapters@set{20261125T010000}{8}% rank-lemma-and-zero-estimate
\chapters@set{20261125T010500}{9}% poincare-series-and-near-orthogonality
\chapters@set{20261124T223000}{10}% elliptic-curves-for-the-proof
\chapters@set{20261125T009500}{11}% fock-model-and-theta-intertwiners
\chapters@set{20261124T223100}{12}% complex-multiplication-and-inert-reduction
\chapters@set{20261125T008000}{13}% one-cm-determinant
\chapters@set{20261124T214500}{14}% supersingular-theta-divisibility
\chapters@set{20261124T214600}{15}% the-isotropic-plane
\chapters@set{20261125T013700}{16}% pairings-on-a-formal-group-of-height-two
\chapters@set{20261125T013500}{17}% hecke-nondegeneracy-on-the-tree
\chapters@set{20261126T030200}{18}% the-bound-is-attained
\chapters@set{20261125T009000}{19}% kronecker-limit-and-chowla-selberg
\chapters@set{20261125T002000}{20}% definite-theta-lifts
\chapters@set{20261125T015500}{21}% siegel-weil-in-the-smallest-case
\chapters@set{20261125T002100}{22}% definite-theta-lifts-and-norms
\chapters@set{20261125T016500}{23}% hecke-operators-and-theta-lifts
\chapters@set{20261125T016700}{24}% divisibility-along-the-chain
\chapters@set{20261125T004000}{25}% denominators-and-sha
\chapters@set{20261125T004100}{26}% tate-twists-and-sha
\chapters@set{20261126T020430}{27}% the-minus-part-by-gauss-sums
\chapters@set{20261125T014100}{28}% congruence-modules
\chapters@set{20261125T014500}{29}% eisenstein-congruences-and-selmer-classes
\chapters@set{20261125T014700}{30}% pseudocharacters-and-the-length-bound
\chapters@set{20261125T011000}{31}% galois-representations-black-box
\chapters@set{20261126T030300}{32}% the-data
\chapters@set{20261126T020210}{33}% bessel-poincare-series-and-their-gram-operators
\chapters@set{20261126T030400}{34}% the-seesaw-identity
\chapters@set{20261126T030500}{35}% divisibility-of-the-composite-lift
\chapters@set{20261126T030600}{36}% theorem-a-short-forms
\chapters@set{20261126T030700}{37}% the-determinant-of-a-chain
\chapters@set{20261126T030800}{38}% the-cost-of-one-lift
\chapters@set{20261126T030900}{39}% denominators-of-projections-to-lifts
\chapters@set{20261126T031000}{40}% twists-with-few-congruences
\chapters@set{20261126T031100}{41}% theorem-c-and-the-main-inequality
\providecommand\chapnumber[1]{%
  \@ifundefined{chapters@#1}%
    {\PackageError{chapters}{Sheet #1 is not listed in the index}%
      {List it in 20261125T012000--grh-workbook__guide.org\MessageBreak
       and run scripts/make-chapters.py.}}%
    {\edef\thechap{\@nameuse{chapters@#1}}}}
\makeatother
\chapnumber{<sheet ID>}
\makeatletter
\providecommand\chapters@set[2]{\@namedef{chapters@#1}{#2}}
\chapters@set{20261125T001000}{1}% zeros-to-ramp-sums
\chapters@set{20261125T001100}{2}% inert-primes-and-estimate-b
\chapters@set{20261125T005000}{3}% galois-fourier-and-averaging
\chapters@set{20261125T005500}{4}% columns-from-a-sphere
\chapters@set{20261125T000000}{5}% slope-method
\chapters@set{20261126T030100}{6}% a-twist-for-each-prime
\chapters@set{20261125T003000}{7}% the-chain-and-the-square-root
\chapters@set{20261125T010000}{8}% rank-lemma-and-zero-estimate
\chapters@set{20261125T010500}{9}% poincare-series-and-near-orthogonality
\chapters@set{20261124T223000}{10}% elliptic-curves-for-the-proof
\chapters@set{20261125T009500}{11}% fock-model-and-theta-intertwiners
\chapters@set{20261124T223100}{12}% complex-multiplication-and-inert-reduction
\chapters@set{20261125T008000}{13}% one-cm-determinant
\chapters@set{20261124T214500}{14}% supersingular-theta-divisibility
\chapters@set{20261124T214600}{15}% the-isotropic-plane
\chapters@set{20261125T013700}{16}% pairings-on-a-formal-group-of-height-two
\chapters@set{20261125T013500}{17}% hecke-nondegeneracy-on-the-tree
\chapters@set{20261126T030200}{18}% the-bound-is-attained
\chapters@set{20261125T009000}{19}% kronecker-limit-and-chowla-selberg
\chapters@set{20261125T002000}{20}% definite-theta-lifts
\chapters@set{20261125T015500}{21}% siegel-weil-in-the-smallest-case
\chapters@set{20261125T002100}{22}% definite-theta-lifts-and-norms
\chapters@set{20261125T016500}{23}% hecke-operators-and-theta-lifts
\chapters@set{20261125T016700}{24}% divisibility-along-the-chain
\chapters@set{20261125T004000}{25}% denominators-and-sha
\chapters@set{20261125T004100}{26}% tate-twists-and-sha
\chapters@set{20261126T020430}{27}% the-minus-part-by-gauss-sums
\chapters@set{20261125T014100}{28}% congruence-modules
\chapters@set{20261125T014500}{29}% eisenstein-congruences-and-selmer-classes
\chapters@set{20261125T014700}{30}% pseudocharacters-and-the-length-bound
\chapters@set{20261125T011000}{31}% galois-representations-black-box
\chapters@set{20261126T030300}{32}% the-data
\chapters@set{20261126T020210}{33}% bessel-poincare-series-and-their-gram-operators
\chapters@set{20261126T030400}{34}% the-seesaw-identity
\chapters@set{20261126T030500}{35}% divisibility-of-the-composite-lift
\chapters@set{20261126T030600}{36}% theorem-a-short-forms
\chapters@set{20261126T030700}{37}% the-determinant-of-a-chain
\chapters@set{20261126T030800}{38}% the-cost-of-one-lift
\chapters@set{20261126T030900}{39}% denominators-of-projections-to-lifts
\chapters@set{20261126T031000}{40}% twists-with-few-congruences
\chapters@set{20261126T031100}{41}% theorem-c-and-the-main-inequality
\providecommand\chapnumber[1]{%
  \@ifundefined{chapters@#1}%
    {\PackageError{chapters}{Sheet #1 is not listed in the index}%
      {List it in 20261125T012000--grh-workbook__guide.org\MessageBreak
       and run scripts/make-chapters.py.}}%
    {\edef\thechap{\@nameuse{chapters@#1}}}}
\makeatother
\chapnumber{<sheet ID>}
\makeatletter
\providecommand\chapters@set[2]{\@namedef{chapters@#1}{#2}}
\chapters@set{20261125T001000}{1}% zeros-to-ramp-sums
\chapters@set{20261125T001100}{2}% inert-primes-and-estimate-b
\chapters@set{20261125T005000}{3}% galois-fourier-and-averaging
\chapters@set{20261125T005500}{4}% columns-from-a-sphere
\chapters@set{20261125T000000}{5}% slope-method
\chapters@set{20261126T030100}{6}% a-twist-for-each-prime
\chapters@set{20261125T003000}{7}% the-chain-and-the-square-root
\chapters@set{20261125T010000}{8}% rank-lemma-and-zero-estimate
\chapters@set{20261125T010500}{9}% poincare-series-and-near-orthogonality
\chapters@set{20261124T223000}{10}% elliptic-curves-for-the-proof
\chapters@set{20261125T009500}{11}% fock-model-and-theta-intertwiners
\chapters@set{20261124T223100}{12}% complex-multiplication-and-inert-reduction
\chapters@set{20261125T008000}{13}% one-cm-determinant
\chapters@set{20261124T214500}{14}% supersingular-theta-divisibility
\chapters@set{20261124T214600}{15}% the-isotropic-plane
\chapters@set{20261125T013700}{16}% pairings-on-a-formal-group-of-height-two
\chapters@set{20261125T013500}{17}% hecke-nondegeneracy-on-the-tree
\chapters@set{20261126T030200}{18}% the-bound-is-attained
\chapters@set{20261125T009000}{19}% kronecker-limit-and-chowla-selberg
\chapters@set{20261125T002000}{20}% definite-theta-lifts
\chapters@set{20261125T015500}{21}% siegel-weil-in-the-smallest-case
\chapters@set{20261125T002100}{22}% definite-theta-lifts-and-norms
\chapters@set{20261125T016500}{23}% hecke-operators-and-theta-lifts
\chapters@set{20261125T016700}{24}% divisibility-along-the-chain
\chapters@set{20261125T004000}{25}% denominators-and-sha
\chapters@set{20261125T004100}{26}% tate-twists-and-sha
\chapters@set{20261126T020430}{27}% the-minus-part-by-gauss-sums
\chapters@set{20261125T014100}{28}% congruence-modules
\chapters@set{20261125T014500}{29}% eisenstein-congruences-and-selmer-classes
\chapters@set{20261125T014700}{30}% pseudocharacters-and-the-length-bound
\chapters@set{20261125T011000}{31}% galois-representations-black-box
\chapters@set{20261126T030300}{32}% the-data
\chapters@set{20261126T020210}{33}% bessel-poincare-series-and-their-gram-operators
\chapters@set{20261126T030400}{34}% the-seesaw-identity
\chapters@set{20261126T030500}{35}% divisibility-of-the-composite-lift
\chapters@set{20261126T030600}{36}% theorem-a-short-forms
\chapters@set{20261126T030700}{37}% the-determinant-of-a-chain
\chapters@set{20261126T030800}{38}% the-cost-of-one-lift
\chapters@set{20261126T030900}{39}% denominators-of-projections-to-lifts
\chapters@set{20261126T031000}{40}% twists-with-few-congruences
\chapters@set{20261126T031100}{41}% theorem-c-and-the-main-inequality
\providecommand\chapnumber[1]{%
  \@ifundefined{chapters@#1}%
    {\PackageError{chapters}{Sheet #1 is not listed in the index}%
      {List it in 20261125T012000--grh-workbook__guide.org\MessageBreak
       and run scripts/make-chapters.py.}}%
    {\edef\thechap{\@nameuse{chapters@#1}}}}
\makeatother
\chapnumber{20261126T040300}
\renewcommand{\thesection}{\thechap.\arabic{section}}

\theoremstyle{definition}
\newtheorem{problem}{Problem}[section]
\theoremstyle{plain}
\newtheorem{theorem}{Theorem}[section]
\newtheorem{corollary}[theorem]{Corollary}

\newcommand{\Z}{\mathbb{Z}}
\newcommand{\Q}{\mathbb{Q}}
\newcommand{\C}{\mathbb{C}}
\newcommand{\Ql}{\bar{\mathbb{Q}}_\ell}
\newcommand{\Zl}{\bar{\mathbb{Z}}_\ell}
\newcommand{\val}{\operatorname{val}}

\setlength{\emergencystretch}{3em}

\title{Roots of unity for many polynomials}
\author{}
\date{November 26, 2026}

\begin{document}
\maketitle
\refstepcounter{chap}\label{chap}

\paragraph{Prerequisites.}
Valuations on \(\Ql\) \cite[Ch.~II]{Ne}; roots of unity of prime power order \cite[Ch.~1--2]{Wa}.  No earlier chapter.

\medskip
Chapter~\ref{tw-chap}, \emph{Twists with few congruences}, describes how the proof bounds Selmer lengths and the valuations of ratios of Hecke \(L\)-values at each prime \(\ell\) through values of polynomials at points \((\zeta_w - 1)_w\), with the \(\zeta_w\) of \(\ell\)-power order.  Suppose finitely many integral polynomials have starting points whose coordinates have valuation at least a fixed rational \(\rho > 0\), and the sum of their starting valuations is bounded.  Can one replace all the points by one tuple \((\zeta_w - 1)_w\) while increasing each valuation by at most a prescribed \(\varepsilon > 0\)?  Theorem~\ref{thm:roots} gives a level depending only on the prime, the bound on the number of variables, \(\rho\), the sum bound and \(\varepsilon\).  Corollary~\ref{pr:powers} treats the first \(m\) powers simultaneously.

Throughout, \(\ell\) is a prime, \(\val\) is the valuation of \(\Ql\) with \(\val(\ell) = 1\) and \(\val(0) = \infty\), \(\lVert y\rVert = \ell^{-\val(y)}\), and \(\Zl = \{y : \val(y) \ge 0\}\).  For \(h \ge 0\) put
\[
  \rho_h := 1/\varphi(\ell^h),
\]
with Euler's function \(\varphi\).  For a polynomial \(g = \sum_i y_i Y^i \in \Ql[Y]\) and a rational \(\rho \ge 0\), the \emph{Gauss valuation} is \(v_\rho(g) := \min_i \bigl(\val(y_i) + \rho i\bigr)\), and \(R_\rho(g)\) is the multiset of roots of \(g\) in \(\Ql\) of valuation exactly \(\rho\), counted with multiplicity; for the zero polynomial \(v_\rho(0) = \infty\) and \(R_\rho(0) = \varnothing\).  Below, \(g\) denotes a nonzero polynomial in \(\Zl[Y]\).  ``Of order \(\ell^h\)'' means of exact order \(\ell^h\).  We write \(\zeta_n\) for a root of unity of exact order \(n\).

\section{One variable}\label{sec:one}

Given \(g\) and a point \(x\) of positive valuation, can one replace \(x\) by \(\zeta - 1\), with \(\zeta\) a root of unity of \(\ell\)-power order, without making \(g\) much more divisible?

\begin{problem}\label{pr:radius}
If \(h \ge 1\) and \(\zeta\) has order \(\ell^h\), then \(\val(\zeta - 1) = \rho_h\).  For \(\ell = 3\) the radii \(\rho_1, \rho_2, \rho_3\) are \(1/2, 1/6, 1/18\); for \(\ell = 2\) they are \(1, 1/2, 1/4\).  The \(\rho_h\) with \(h \ge 1\) are distinct and tend to \(0\).
\end{problem}

\begin{problem}\label{pr:gauss}
Let \(0 \ne g \in \Zl[Y]\), and write \(g = c\prod_k (Y - \alpha_k)\) in \(\Ql[Y]\).  For rational \(\rho \ge 0\),
\[
  v_\rho(g) = \val(c) + \sum\nolimits_k \min\bigl(\val(\alpha_k), \rho\bigr).
\]
Hence \(v_\rho(g_1g_2) = v_\rho(g_1) + v_\rho(g_2)\), and \(\rho \mapsto v_\rho(g)\) is nondecreasing.
\end{problem}

\begin{problem}\label{pr:at-point}
\textup{(a)} If \(\val(x) \ge \rho\), then \(\val(g(x)) \ge v_\rho(g)\).
\textup{(b)} If \(\val(y) = \rho\), then
\[
  \val(g(y)) = v_\rho(g) + \sum\nolimits_{\alpha \in R_\rho(g)} \bigl(\val(y - \alpha) - \rho\bigr).
\]
\end{problem}

\begin{problem}\label{pr:example}
Let \(\ell = 3\) and \(g = Y^2 + 3Y + 3\).  Then \(v_\rho(g) = \min(1, 2\rho)\), and for \(\zeta\) of order \(3^h\) the valuation \(\val(g(\zeta - 1))\) is \(\infty\) if \(h = 1\) and \(2\rho_h = v_{\rho_h}(g)\) if \(h \ge 2\).
\end{problem}

We bound the root counts on the radii \(\rho_h \le \rho\), each weighted by \(\rho_h\).  For a rational \(\rho > 0\) and an integer \(h_0 \ge 1\), let \(H(\rho, h_0)\) be the set of \(h\) with \(1 \le h \le h_0\) and \(\rho_h \le \rho\).

\begin{problem}\label{pr:few-radii}
\(\sum_{h \in H(\rho, h_0)} \# R_{\rho_h}(g) \cdot \rho_h \le v_\rho(g)\).
\end{problem}

\begin{problem}\label{pr:pigeon}
Let \(\delta > 0\) and \(\beta \ge 0\) (in this sheet \(\beta\) is a real bound), with \(\beta < \delta \sum_{h \in H(\rho, h_0)} 1/(h + 1)\).  If \(v_\rho(g) \le \beta\), then some \(h \in H(\rho, h_0)\) has \(\# R_{\rho_h}(g) \cdot \rho_h \le \delta/(h + 1)\).
\end{problem}

Next, how close can the \(\zeta - 1\) of a given order be to a fixed point?

\begin{problem}\label{pr:close-to-few}
Let \(\alpha \in \Zl\) with \(\val(\alpha) > 0\), and \(h \ge 1\).  Then
\[
  \sum\nolimits_{\zeta} \min\bigl(\val(\zeta - 1 - \alpha), 1\bigr) \le h + 1,
\]
the sum over the \(\zeta\) of order \(\ell^h\).  For \(\ell = 3\), \(h = 1\) and \(\alpha = \zeta_0 - 1\), with \(\zeta_0\) of order \(3\), the sum is \(3/2\).
\end{problem}

\begin{problem}\label{pr:ninth}
For \(\ell = 3\), \(h = 2\) and \(\alpha = \zeta_9 - 1\), the six terms \(\min(\val(\zeta - 1 - \alpha), 1)\), for \(\zeta = \zeta_9^e\) with \(e = 1, 2, 4, 5, 7, 8\), are \(1, \tfrac16, \tfrac12, \tfrac16, \tfrac12, \tfrac16\); their sum is \(\tfrac52 \le h + 1\).
\end{problem}

\begin{problem}\label{pr:average}
Let \(R\) be a finite multiset of elements of \(\Zl\) of positive valuation, and \(h \ge 1\).  Some \(\zeta\) of order \(\ell^h\) has \(\sum_{\alpha \in R} \min\bigl(\val(\zeta - 1 - \alpha), 1\bigr) \le \#R \cdot \rho_h (h + 1)\).
\end{problem}

\begin{problem}\label{pr:one-variable}
Let \(0 < \delta < 1\), \(\rho > 0\) rational, and \(\beta < \delta\sum_{h \in H(\rho, h_0)} 1/(h + 1)\).  If \(v_\rho(g) \le \beta\), then for some \(h \in H(\rho, h_0)\) and some \(\zeta\) of order \(\ell^h\),
\[
  \val\bigl(g(\zeta - 1)\bigr) \le v_{\rho_h}(g) + \delta .
\]
\end{problem}

\begin{problem}\label{pr:harmonic}
For rational \(\rho > 0\), real \(\delta > 0\) and \(\beta \ge 0\), some \(h_0\) has \(\beta < \delta\sum_{h \in H(\rho, h_0)} 1/(h + 1)\).  For \(\ell = 2\), \(\rho = \delta = 1/2\) and \(\beta = 1\), the least such \(h_0\) is \(18\).
\end{problem}

\begin{problem}\label{pr:one-var-point}
Let \(\delta\), \(\rho\), \(\beta\), \(h_0\) be as in Problem~\ref{pr:one-variable}.  If \(\val(x) \ge \rho\) and \(\val(g(x)) \le \beta\), then some \(\zeta\) with \(\zeta^{\ell^{h_0}} = 1\) and \(\val(\zeta - 1) \le \rho\) has \(\val(g(\zeta - 1)) \le \val(g(x)) + \delta\).
\end{problem}

\section{Several variables and several powers}\label{sec:several}

Let \(\rho > 0\) be rational, and \(\mathcal V\) and \(I\) finite sets.  For \(i \in I\) let \(F_i \in \Zl[X_j : j \in \mathcal V]\), and let \(x_i \in \Zl^\mathcal V\) be a point.  We seek one tuple of roots of unity that works for every \(i\).

\begin{problem}\label{pr:one-coordinate}
Fix \(j \in \mathcal V\), and assume \(\val(x_{ij}) \ge \rho\) for all \(i\) (the other coordinates are only integral).  Let \(0 < \delta < 1\), \(\beta < \delta\sum_{h \in H(\rho, h_0)} 1/(h + 1)\), and \(\sum_i \val(F_i(x_i)) \le \beta\).  There is \(\zeta\) with \(\zeta^{\ell^{h_0}} = 1\) and \(\val(\zeta - 1) \le \rho\) such that, if the coordinate \(x_{ij}\) is replaced by \(\zeta - 1\) for every \(i\), each \(\val(F_i(x_i))\), and also their sum, grows by at most \(\delta\).
\end{problem}

\begin{problem}\label{pr:all-coordinates}
Let \(0 < \delta < 1\), \(\#\mathcal V \cdot \delta \le 1\), \(\beta + 1 < \delta\sum_{h \in H(\rho, h_0)} 1/(h + 1)\), \(\val(x_{ij}) \ge \rho\) for all \(i, j\), and \(\sum_i \val(F_i(x_i)) \le \beta\).  There are \(\zeta_j\), \(j \in \mathcal V\), with \(\zeta_j^{\ell^{h_0}} = 1\), \(\val(\zeta_j - 1) \le \rho\), and \(\val\bigl(F_i((\zeta_j - 1)_j)\bigr) \le \val(F_i(x_i)) + \#\mathcal V\cdot\delta\) for all \(i\).
\end{problem}

\begin{problem}\label{pr:two-coordinates}
Let \(\ell = 3\), \(F_1 = X + Y\), \(F_2 = XY + 3\), \(F_3 = X^2 + 3X + 3\), with starting points \((3, 0)\), \((0, 3)\), \((0, 0)\), at which each \(F_i\) has valuation \(1\).  For \(\zeta = \zeta_{27}\), the common point \((\zeta - 1, \zeta^2 - 1)\) gives \(\val(F_i) = \tfrac19\) for \(i = 1, 2, 3\).
\end{problem}

The following theorem is proved by the problems of this section.  In Theorem~\ref{thm:roots}, the level \(h_0\) depends only on \(\ell, s_0, \rho, \beta, \varepsilon\).

\begin{theorem}\label{thm:roots}
Let \(s_0 \ge 0\) be an integer, \(\rho > 0\) rational, \(\beta \ge 0\) and \(\varepsilon > 0\).  There is \(h_0\) with the following property.  Let \(\mathcal V\), \(I\) be finite sets with \(\#\mathcal V \le s_0\), let \(F_i \in \Zl[X_j : j \in \mathcal V]\) and \(x_i \in \Zl^\mathcal V\) for \(i \in I\), with \(\val(x_{ij}) \ge \rho\) for all \(i, j\) and \(\sum_i \val(F_i(x_i)) \le \beta\).  Then there are \(\zeta_j \in \Zl\), \(j \in \mathcal V\), with
\[
  \zeta_j^{\ell^{h_0}} = 1,\quad \val(\zeta_j-1)\le\rho
     \quad(j\in\mathcal V),\qquad
  \val\bigl(F_i((\zeta_j-1)_j)\bigr)
     \le \val\bigl(F_i(x_i)\bigr)+\varepsilon \quad(i\in I).
\]
\end{theorem}

\medskip
A twist will be used through its first \(m\) powers at once, so the roots of unity must also serve for \(\zeta_j^n - 1\), \(1 \le n \le m\).

\begin{problem}\label{pr:root-of-x}
Let \(n \ge 1\) and \(x \in \Zl\) with \(\val(x) \ge \rho\).  Some \(x' \in \Zl\) has \((1 + x')^n = 1 + x\) and \(\val(x') \ge \rho/n\).
\end{problem}

\begin{problem}\label{pr:powers-far}
Let \(0 < \rho < 1\), \(1 \le n \le m\), and let \(\zeta\) have \(\ell\)-power order and \(\val(\zeta - 1) \le \rho/m\).  Then \(\val(\zeta^n - 1) \le \rho\); in particular \(\zeta^n \ne 1\).  For \(\ell = 2\), \(\zeta\) of order \(8\), \(m = n = 2\) and \(\rho = 1/2\), equality holds.
\end{problem}

\begin{problem}\label{pr:powers-table}
Let \(\ell = 3\) and \(\Phi_3(Y) = Y^2 + Y + 1\).  The values \(\val(\Phi_3(\zeta_n^j))\) for \(j = 1, 2, 3\) are \(\tfrac13, \tfrac13, \infty\) for \(n = 9\); \(\tfrac19, \tfrac19, \tfrac13\) for \(n = 27\); and \(\tfrac1{27}, \tfrac1{27}, \tfrac19\) for \(n = 81\).  For \(m = 3\), \(\rho = 1/2\), one variable, \(F(X) = \Phi_3(1 + X)\), \(x = 0\) and \(\beta = 1\), no root of order \(9\) satisfies the conclusion of Corollary~\ref{pr:powers}, for any \(\varepsilon > 0\).
\end{problem}

Problems~\ref{pr:root-of-x} and~\ref{pr:powers-far} and Theorem~\ref{thm:roots} give the form used in Chapter~\ref{tw-chap}:

\begin{corollary}\label{pr:powers}
Let integers \(s_0 \ge 0\), \(m \ge 1\), a rational \(\rho > 0\), \(\varepsilon > 0\), a finite set \(I\) and numbers \(\beta_i \ge 0\), \(i \in I\), be given.  There is \(h_0\) such that, for every finite set \(\mathcal V\) with \(\#\mathcal V \le s_0\), all \(F_i \in \Zl[X_j : j \in \mathcal V]\) and all \(x_i \in \Zl^\mathcal V\) satisfying \(\val(x_{ij}) \ge \rho\) for every \(i, j\) and \(\val(F_i(x_i)) \le \beta_i\) for every \(i\), some \(\zeta_j\) with \(\zeta_j^{\ell^{h_0}} = 1\) and \(\val(\zeta_j - 1) \le \rho/m\) satisfy
\[
  \val\bigl(F_i((\zeta_j^n - 1)_j)\bigr) \le \val\bigl(F_i(x_i)\bigr) + \varepsilon \qquad (i \in I,\ 1 \le n \le m).
\]
\end{corollary}

\medskip
The exercises derive Theorem~\ref{thm:roots} and Corollary~\ref{pr:powers} from the stated prerequisites and the divergence of the harmonic series.  Chapter~\ref{w2-chap} states the polynomial input, and Chapter~\ref{tw-chap} uses Corollary~\ref{pr:powers} to construct the twist and bound its powers.

\paragraph{In the Lean formalization.}
In the Lean proof \cite{GL}, Theorem~\ref{thm:roots} is \texttt{New.Selmer.exists\_rootsOfUnity\_ellVal\_eval\_le\_add} (\nolinkurl{New/Congruences/Selmer/RootsOfUnity.lean}); two one-variable steps, used in Problems~\ref{pr:pigeon} and~\ref{pr:one-variable}, are \texttt{New.Selmer.exists\_card\_rootsOfValuation\_mul\_le} and \texttt{New.Selmer.exists\_rootOfUnity\_eval\_le\_gaussVal\_add} in the same file; the form for powers, Corollary~\ref{pr:powers}, is \texttt{New.Selmer.exists\_rootsOfUnity\_ellVal\_eval\_le} (\nolinkurl{New/Congruences/Selmer/GeneralTheorem.lean}).  The truncated valuation \(\max(\val, 0)\) is \nolinkurl{New.ellVal}.  The polynomial input and the use of the first \(m\) powers of one twist are \nolinkurl{New.Selmer.exists_twistPolynomials} (\nolinkurl{New/Congruences/Selmer/Line.lean}) and the proof of \nolinkurl{New.Selmer.exists_anticyclotomic_twist} (\nolinkurl{New/Congruences/Selmer/GeneralTheorem.lean}).  The passage from a Stickelberger value to a bound on a Hecke \(L\)-value ratio uses \nolinkurl{New.Selmer.exists_stickelbergerSystem_ellVal_ratio_le} (\nolinkurl{New/Congruences/Selmer/LValueSystems.lean}).

\section*{Hints}

\subsection*{Section~\ref{sec:one}}
\begin{description}
\item[\ref{pr:radius}] The conjugates of \(\zeta - 1\) have product \(\pm\Phi_{\ell^h}(1) = \pm\ell\).
\item[\ref{pr:gauss}] For \(\rho = 0\) reduce modulo the maximal ideal of \(\Zl\); for \(\rho = \val(\pi)\) substitute \(Y = \pi Z\).
\item[\ref{pr:at-point}] For (b), factor \(g\) and use \(\val(y - \alpha) = \min(\val(\alpha), \rho)\) when \(\val(\alpha) \ne \rho\).
\item[\ref{pr:example}] \(g(Y) = ((1 + Y)^3 - 1)/Y\).
\item[\ref{pr:few-radii}] By Problem~\ref{pr:radius} a root lies in at most one \(R_{\rho_h}\); compare with Problem~\ref{pr:gauss}, using \(v_0(g) \ge 0\).
\item[\ref{pr:close-to-few}] The \(\zeta - 1 - \alpha\) are roots of \((1 + \alpha + Y)^{\ell^h} - 1\); look at its coefficient of \(Y\) and use Problem~\ref{pr:gauss} with \(\rho = 1\).
\item[\ref{pr:ninth}] \(\zeta_9^e - \zeta_9 = \zeta_9(\zeta_9^{e-1} - 1)\); Problem~\ref{pr:radius}.
\item[\ref{pr:average}] Sum Problem~\ref{pr:close-to-few} over \(R\) and average over the \(1/\rho_h\) choices of \(\zeta\).
\item[\ref{pr:one-variable}] Choose the radius by Problem~\ref{pr:pigeon} and \(\zeta\) by Problem~\ref{pr:average}.
\item[\ref{pr:harmonic}] The harmonic series diverges.
\item[\ref{pr:one-var-point}] Problem~\ref{pr:at-point}(a) and the monotonicity in Problem~\ref{pr:gauss}.
\end{description}

\subsection*{Section~\ref{sec:several}}
\begin{description}
\item[\ref{pr:one-coordinate}] Apply Problem~\ref{pr:one-variable} to the product of the one-variable polynomials \(g_i(Y)\) obtained by freeing the coordinate \(j\); for each \(i\) separately, use \(\val(g_i(\zeta - 1)) \ge v_{\rho_h}(g_i)\).
\item[\ref{pr:all-coordinates}] Induction on the number of coordinates moved, with Problem~\ref{pr:one-coordinate}, keeping the running sum at most \(\beta + 1\).  For Theorem~\ref{thm:roots} take \(\delta = \min(\varepsilon, 1)/(s_0 + 2)\) and \(h_0\) from Problem~\ref{pr:harmonic} with \(\beta + 1\) in place of \(\beta\).
\item[\ref{pr:two-coordinates}] Factor \(F_1\) and \(F_3\) at the point; Problem~\ref{pr:radius}.
\item[\ref{pr:root-of-x}] The roots of \((1 + Y)^n - 1 - x\) have product \(\pm x\).
\item[\ref{pr:powers-far}] First compare \(m\) with \(\varphi(\ell^h)\), where \(\ell^h\) is the order of \(\zeta\).  What is the order of \(\zeta^n\)?
\item[\ref{pr:powers-table}] \(\Phi_3(y) = (y^3 - 1)/(y - 1)\); Problem~\ref{pr:radius}.
\item[Corollary~\ref{pr:powers}] Theorem~\ref{thm:roots} with radius \(\rho/m\), after the substitution \(X_j \mapsto (1 + X_j)^n - 1\); Problem~\ref{pr:root-of-x}.
\end{description}

\begin{thebibliography}{GL}
\bibitem[GL]{GL} L.~Sernau, \emph{GRH for Dirichlet \(L\)-functions, in Lean~4}, a shorter rewriting of the Lean proof, text by Claude, \url{https://github.com/lsernau/grh-lean}, commit \texttt{72fead1}, readable edition \texttt{grh/}.
\bibitem[Ne]{Ne} J.~Neukirch, \emph{Algebraic Number Theory}, Grundlehren Math.\ Wiss.\ 322, Springer, 1999.
\bibitem[Wa]{Wa} L.~C.~Washington, \emph{Introduction to Cyclotomic Fields}, 2nd ed., Grad.\ Texts in Math.\ 83, Springer, 1997.
\end{thebibliography}

\end{document}
